\section{Discussion}
\markright{Discussion}

The forward-bias readings show the expected nonlinear behavior of the silicon p-n junction
diode, with current increasing steadily as the applied voltage rises and the depletion barrier
is reduced. However, a discrepancy was observed between the theoretical cut-in voltage of
0.7 V and the experimental results, where the diode voltage (VF) only reached a maximum
of 0.48V. This is attributable to the 10kΩ series resistor being too large, which limited the
forward current to a range where the diode remains in the ”knee” region of its characteristic
curve. In reverse bias, the current remained extremely small; with the 10kΩ resistor, our
multimeter did not display measurable reverse current, which is consistent with the diode
blocking majority carrier flow. If a 1kΩ resistor had been used, the full 0.7V conduction
threshold and a larger reverse leakage current would likely have been moreobservable. Small
differences between simulation and observation are attributable to real-diode effects such as
tolerance, meter loading, and internal resistance.
